\section{An incoming-orientation invariant for six existing increment entries}
\label{sec:increment-orientation}

The original grammar restricts a terminal to the incoming bit values observed
in the source from which the grammar was extracted. An observed restriction
need not be a semantic restriction. Three previously one-sided terminals can
be made available on both sides without changing any binary transition.

Use the old state labels. Let $B(v)=0\,1^{v+1}$ and let $B$ concatenate finite
lists. At a dispatcher frame, let the pointer word $p$ occupy cells
$0,\ldots,\ell-1$, with marker one at $-1$, header one at $D$, and complete
register suffix $B(R)B(S)$. Here $R\in\mathbb N^9$, $S$ is a nonempty finite
list, $D-\ell\geq3$, and $p$ contains no seven consecutive zeros. The
canonical Routing245 bounds $\ell\leq51$ and $D\geq64$ imply the stated gap
condition. Translation of every coordinate is harmless.

Assume that the transparent decoder follows the pointer $p$, preserving it,
and enters an atomic state $Q$ at cell $\ell$. On each internal return to
the dispatcher, the same unchanged pointer follows the same decoder route.
The following existing entries have rows
\[
 Q/0=0R57,\qquad Q/1=0R\operatorname{skip}_r.
\]
\begin{center}
\begin{tabular}{c|c|c|c}
$Q$&logical increment&physical register $r$&$\operatorname{skip}_r$\\\hline
39&$X$&8&69\\
70&$T$&2&31\\
72&$H$&7&64\\
93&$Y$&5&40\\
169&$D$&3&56\\
201&$A$&4&38
\end{tabular}
\end{center}
The last two entries already exist in the framework but are not singleton
tokens in the present thirty-token fine alphabet.

\paragraph{Orientation lemma.}
Each displayed entry implements its indicated increment for both values of
the last bit of $p$. It increases exactly register $r$, appends a spare of
value one, and performs the full-width binary successor of the pointer,
including its marker shift in the overflow case. No extra state is required.

\paragraph{Proof.}
The proof decomposes the literal execution into four phases. The only
register assumptions used are the already checked mark and increment
contracts, valid for arbitrary natural-valued registers and arbitrary
nonempty spare suffixes. Those contracts use identical rows in Routing245.

First, the rows $Q/0$, $57/0$ enter state 42 at cell $\ell$. Its zero
self-loop fills the entire gap with ones. At the header, $42/1$ returns
left through state 2 and the marker scanner. The scanner detects eight
successive zeros, while the pointer has at most six. It therefore returns
to the left marker without changing the tape. The pointer is unchanged;
the decoder returns to $Q$, now reading one at cell $\ell$.

Second, $Q/1$ clears that cell and enters the appropriate selector to its
right. The selector crosses the remaining gap ones and header, marks
register $r$, and returns left from the register file. Both possible mark
return states, 2 and 6, have the same row $1L4$ on the header one. The
scanner consequently sees the original pointer followed by one extra zero
at cell $\ell$. Its zero runs have length at most seven, still below the
eight-zero detector threshold. It returns to the unchanged pointer, and
the transparent decoder again enters $Q$ at cell $\ell$.

Third, $Q/0$ moves to state 57 at cell $\ell+1$, whose one is cleared by
$57/1$. State 73 starts at $\ell+2$, crosses the remaining gap and header,
and invokes the marked-register increment contract. That contract increases
only register $r$, appends $B(1)$, and returns in state 48 to the left of
the register file. The rows of 48 and 63 cross the gap ones until the two
cleared cells are reached. In particular, $63/0$ puts state 48 at cell
$\ell$, and $48/0$ puts state 7 on the final pointer bit at $\ell-1$.
No marker scan occurs across the potentially longer zero run formed by the
pointer and these two temporary zero cells.

Finally, state 7 performs binary carry to the left: ones become zeros until
the first zero is changed to one. If every bit of $p$ was one, it also
clears the old marker and writes the new marker one cell farther left.
States 10 and 14 clear the remaining temporary gap ones. State 20 restores
the header one and returns through state 2. State 2 skips the complete
trailing zero suffix before invoking the marker detector, so the long zero
suffix produced by carry is not mistaken for a marker gap. Every other
zero run in the surviving pointer prefix has length at most six. The final
dispatcher frame therefore has precisely the full-width successor pointer,
the indicated register increment, the new spare, and the prescribed
overflow displacement. Every phase takes a finite positive number of
steps; the register phases are justified by unary and finite-list induction.

None of these phases branches on the incoming bit as a special case. The
only read of that bit that matters is the ordinary binary carry in the final
phase, which correctly handles both zero and one. This proves the lemma.

\paragraph{Consequences and limits.}
The previously restricted original symbols $h,m,r$ (states 39, 72, 93)
may use orientation mask three. Their fine tokens are respectively
$A,T,F$. State 70 already had both orientations. State 18, implementing
increment of $N$, has a different zero row and is outside the lemma: a
concrete outer pointer $00$ fails its opposite-orientation contract.
Broader atomic experiments in this package quantify register values and
header slack but use finitely many concrete outer prefixes; those
experiments are not generalized by this lemma. A complete new candidate
must still replay all its actual terminal paths on its own literal table.
